El protocolo \textbf{Open Shortest Path First (OSPF)} fue desarrollado por el IETF como una alternativa robusta y escalable a los protocolos de vector-distancia. Es el estándar de facto para el enrutamiento interior en redes corporativas y de proveedores de servicios.

\subsection{Algoritmo de Estado de Enlace y Dijkstra}

A diferencia de RIP, OSPF no intercambia tablas de rutas. En su lugar, utiliza un enfoque de \textbf{estado de enlace}:

\begin{enumerate}
    \item \textbf{Base de Datos de Topología:} Cada enrutador mantiene una base de datos (\textit{Link State Database} o LSDB) que representa el mapa completo de la red.
    \item \textbf{Inundación de LSA:} Cuando ocurre un cambio, los enrutadores difunden anuncios de estado de enlace (\textit{Link State Advertisements}) a toda la red para sincronizar las LSDB.
    \item \textbf{Algoritmo SPF:} Con la LSDB actualizada, cada enrutador ejecuta localmente el \textbf{algoritmo de Dijkstra} para construir un árbol de caminos mínimos, situándose a sí mismo como la raíz, para determinar la mejor ruta hacia cada destino.
\end{enumerate}

\subsection{Características Avanzadas}

OSPF ofrece funcionalidades que mejoran la eficiencia y la gestión:

\begin{itemize}
    \item \textbf{Jerarquía de Áreas:} Permite dividir un Sistema Autónomo en áreas (siendo el Área 0 el núcleo o \textit{backbone}). Esto limita la propagación de LSA y reduce el tamaño de las LSDB en redes extensas.
    \item \textbf{Métrica de Costo:} Utiliza el "costo" como métrica, el cual suele basarse en el ancho de banda del enlace. Esto permite una selección de rutas mucho más precisa que el simple conteo de saltos.
    \item \textbf{Balanceo de Carga:} Soporta el envío de tráfico a través de múltiples rutas de igual costo (\textit{Equal-Cost Multi-Path}).
    \item \textbf{Autenticación:} Permite asegurar los intercambios de mensajes para prevenir la inyección de rutas falsas.
\end{itemize}

\subsection{OSPFv2 y OSPFv3}

\begin{itemize}
    \item \textbf{OSPFv2 (IPv4):} Los mensajes se encapsulan directamente en IP (protocolo 89). Utiliza un encabezado fijo de 24 octetos y gestiona la autenticación de forma nativa.
    \item \textbf{OSPFv3 (IPv6):} Representa un rediseño para lograr \textbf{independencia del protocolo de red}. Los identificadores (Router ID y Area ID) siguen siendo de 32 bits, desligándose de las direcciones IPv6. Además, delega la seguridad al encabezado de autenticación de IPv6 (IPsec), eliminando los campos de seguridad del protocolo base.
\end{itemize}
